% Part: model-theory
% Chapter: models-of-arithmetic
% Section: computable-models

\documentclass[../../../include/open-logic-section]{subfiles}

\begin{document}

\olfileid{mod}{mar}{cmp}
\section{અંકગણિતના સંગણનીય નિદર્શો}

\begin{explain}
પ્રમાણભૂત નિદર્શ~$\Struct{N}$નાં બે સરસ લક્ષણો છે. તેનું
વ્યાપકક્ષેત્ર પ્રાકૃતિક સંખ્યાઓ~$\Nat$ છે, એટલે કે તેના ઘટકો એવી જ
વસ્તુઓ છે જેના વિશે અંકગણિતની ભાષાથી વાત કરવી છે, અને પ્રમાણભૂત આંકિક
પદ~$\num{n}$ ખરેખર~$n$ને જ નિર્દેશે છે. બીજું સરસ લક્ષણ એ છે કે
$\Lang{L_A}$ના અતાર્કિક સંકેતોનાં બધાં અર્થઘટનો \emph{સંગણનીય} છે.
$\Assign{\prime}{N}$, $\Assign{+}{N}$ અને $\Assign{\times}{N}$ તરીકે
કામ કરતાં ઉત્તરગામી, સરવાળા અને ગુણાકારનાં વિધેયો સંખ્યાઓનાં સંગણનીય
વિધેયો છે. (માનો કે ત્યૂરિંગ યંત્રો વડે સંગણનીય, અથવા આદિમ
પુનરાવર્તનથી વ્યાખ્યેય.) અને $\Struct{N}$ પરનો લઘુ-કરતાં સંબંધ, એટલે
$\Assign{<}{N}$, નિર્ણેય છે.

$\Th{Q}$ અને $\Th{PA}$ જેવા અંકગણિતીય સિદ્ધાંતોના અપ્રમાણભૂત
નિદર્શોમાં અપ્રમાણભૂત ઘટકો હોવા જ જોઈએ. તેથી તેમનાં વ્યાપકક્ષેત્રોમાં
સામાન્ય રીતે $\Nat$ ઉપરાંત બીજા !!{element}s પણ હોય છે. તેમ છતાં, કોઈ
પણ ગણનીય !!{structure}ને કોઈ પણ !!{denumerable} ગણ પર, જેમાં $\Nat$
પણ સામેલ છે, બનાવી શકાય છે. તેથી વ્યાપકક્ષેત્ર~$\Nat$ ધરાવતા
અપ્રમાણભૂત નિદર્શો પણ છે. આવા નિદર્શ~$\Struct{M}$માં ઓછામાં ઓછી કેટલીક
સંખ્યાઓ પોતાની સામાન્ય ભૂમિકા ભજવી શકે નહિ, કારણ કે એવો કોઈ~$k$ હોવો
જોઈએ જે દરેક $n\in\Nat$ માટે $\Value{\num{n}}{M}$થી જુદો હોય.
\end{explain}

\begin{defn}
$\Lang{L_A}$ માટેની !!^a{structure}~$\Struct{M}$ \emph{સંગણનીય}
ત્યારે અને માત્ર ત્યારે જ્યારે $\Domain{M}=\Nat$,
$\Assign{\prime}{M}$, $\Assign{+}{M}$ અને $\Assign{\times}{M}$
સંગણનીય વિધેયો હોય અને $\Assign{<}{M}$ નિર્ણેય સંબંધ હોય.
\end{defn}

\begin{ex}\ollabel{ex:comp-model-q}
\olref[mdq]{ex:model-K-of-Q}ની સંરચના~$\Struct{K}$ યાદ કરો. તેનું
વ્યાપકક્ષેત્ર $\Domain{K}=\Nat\cup\{a\}$ અને અર્થઘટનો આ હતાં:
\begin{align*}
  \Assign{\Obj{0}}{K} & = 0\\
  \Assign{\prime}{K}(x) & =
  \begin{cases}
    x+1 & \text{જો $x\in\Nat$}\\
    a & \text{જો $x=a$}
  \end{cases}\\
  \Assign{+}{K}(x,y) & =
  \begin{cases}
    x+y & \text{જો $x$, $y\in\Nat$}\\
    a & \text{અન્યથા}
  \end{cases}\\
  \Assign{\times}{K}(x,y) & =
  \begin{cases}
    xy & \text{જો $x$, $y\in\Nat$}\\
    0 & \text{જો $x=0$ અથવા $y=0$}\\
    a & \text{અન્યથા}\\
  \end{cases}\\
  \Assign{<}{K} & =
  \Setabs{\tuple{x,y}}{x,y\in\Nat\text{ અને }x<y}\cup
  \Setabs{\tuple{x,a}}{x\in\Domain{K}}
\end{align*}
\sourcecorrection{OLMAR-014}{સ્થિર અંગ્રેજી મૂળની
  \texttt{computable-models.tex}, પંક્તિ~66માં બીજા ગણ-નિર્માણનું
  ક્રમયુક્ત જોડ $\tuple{x,a}$ છે, પરંતુ તેની શરતમાં અજોડ ચલ~$n$ વપરાય
  છે. અહીં \olref[mdq]{ex:model-K-of-Q}ની મૂળ વ્યાખ્યા મુજબ
  $x\in\Domain{K}$ પુનઃસ્થાપિત કર્યું છે.}
પરંતુ $\Domain{K}$ !!{denumerable} છે અને તેથી $\Nat$ જેટલી જ સંખ્યા
ધરાવે છે. દાખલા તરીકે, $g(0)=a$ અને $n>0$ માટે $g(n)=n-1$ ધરાવતું
$g\colon\Nat\to\Domain{K}$ !!a{bijection} છે. તેને $\Th{Q}$ના નવા
નિદર્શ~$\Struct{K'}$ અને $\Struct{K}$ વચ્ચેની એકરૂપતા બનાવી શકીએ છીએ.
$\Struct{K'}$માં $\Lang{L_A}$ના સંકેતોને જુદાં વિધેયો અને સંબંધો
આપવા પડે, કારણ કે $\Nat$ના જુદા !!{element}s પ્રમાણભૂત અને અપ્રમાણભૂત
સંખ્યાઓની ભૂમિકાઓ ભજવે છે.
\sourcecorrection{OLMAR-015}{સ્થિર અંગ્રેજી મૂળની
  \texttt{computable-models.tex}, પંક્તિઓ~69--70માં $g(0)=a$ અને
  $n>0$ માટે $g(n)=n+1$ ધરાવતું વિધેય દ્વિએક-વ્યાપ્ત કહેવાયું છે, પરંતુ
  તે $0$ અને~$1$ને પરાસમાંથી છોડે છે. તરત પછીની રચનામાં પ્રમાણભૂત
  $0$ની ભૂમિકા~$1$ ભજવે છે; તેના અને આપેલાં બધા પરિવહિત અર્થઘટનો સાથે
  સુસંગત સાચું વિધેય $g(n)=n-1$ છે.}

ખાસ કરીને, હવે $0$, લઘુતમ પ્રમાણભૂત સંખ્યાની નહિ પરંતુ~$a$ની ભૂમિકા
ભજવે છે. હવે લઘુતમ પ્રમાણભૂત સંખ્યા~$1$ છે. તેથી
$\Assign{\Obj{0}}{K'}=1$ રાખીએ. ઉત્તરગામી વિધેય પણ હવે જુદું છે:
પ્રમાણભૂત સંખ્યા, એટલે $n>0$, આપતાં તે હજી $n+1$ આપે છે. પરંતુ હવે
$0$, પોતાનો જ ઉત્તરગામી એવા~$a$ની ભૂમિકા ભજવે છે. તેથી
$\Assign{\prime}{K'}(0)=0$. સરવાળા અને ગુણાકાર માટે એવી જ રીતે
\begin{align*}
\Assign{+}{K'}(x,y) & =
  \begin{cases}
    x+y-1 & \text{જો $x$, $y>0$}\\
    0 & \text{અન્યથા}
  \end{cases}\\
\Assign{\times}{K'}(x,y) & =
  \begin{cases}
    1 & \text{જો $x=1$ અથવા $y=1$}\\
    xy-x-y+2 & \text{જો $x$, $y>1$}\\
    0 & \text{અન્યથા}\\
  \end{cases}
\end{align*}
રાખીએ. અને $\tuple{x,y}\in\Assign{<}{K'}$ ત્યારે અને માત્ર ત્યારે
જ્યારે $x<y$, $x>0$ અને $y>0$ હોય, અથવા $y=0$ હોય.

આ બધાં વિધેયો પ્રાકૃતિક સંખ્યાઓનાં સંગણનીય વિધેયો છે અને
$\Assign{<}{K'}$, $\Nat$ પરનો નિર્ણેય સંબંધ છે—પરંતુ તેઓ $\Nat$ પરના
ઉત્તરગામી, સરવાળા અને ગુણાકારનાં વિધેયો જેવા નથી, અને
$\Assign{<}{K'}$, $\Nat$ પરના $<$ સંબંધ જેવો નથી.
\end{ex}

\begin{prob}
$\Domain{L'}=\Nat$ ધરાવતી અને
\olref[mod][mar][mdq]{ex:model-L-of-Q}ની $\Struct{L}$ને એકરૂપ એવી
!!a{structure}~$\Struct{L'}$ આપો.
\end{prob}

\begin{explain}
\Olref{ex:comp-model-q} બતાવે છે કે $\Th{Q}$ને વ્યાપકક્ષેત્ર~$\Nat$
ધરાવતા સંગણનીય અપ્રમાણભૂત નિદર્શો છે. પરંતુ નીચેનું પરિણામ બતાવે છે
કે $\Th{PA}$ના (અને તેથી $\Th{TA}$ના પણ) નિદર્શો માટે આ સત્ય નથી.
\end{explain}

\begin{thm}[ટેનેનબાઉમનું પ્રમેય]
$\Th{PA}$નો દરેક સંગણનીય નિદર્શ પ્રમાણભૂત છે, એટલે કે
$\Struct{N}$ને એકરૂપ છે.
\sourcecorrection{OLMAR-016}{સ્થિર અંગ્રેજી મૂળની
  \texttt{computable-models.tex}, પંક્તિ~120 શાબ્દિક રીતે
  $\Struct{N}$ને એકમાત્ર સંગણનીય નિદર્શ કહે છે. પરંતુ $\Nat$ની કોઈ
  સંગણનીય આપલે વડે તેનાં અર્થઘટનો પરિવહિત કરવાથી એ જ ક્ષેત્ર પર જુદી
  સંગણનીય સંરચનાઓ મળે છે. ટેનેનબાઉમના પ્રમેયનો ચોક્કસ નિષ્કર્ષ એ છે
  કે કોઈ સંગણનીય \emph{અપ્રમાણભૂત} નિદર્શ નથી; તેથી અહીં દરેક સંગણનીય
  નિદર્શ $\Struct{N}$ને એકરૂપ છે એમ લખ્યું છે.}
\end{thm}

\end{document}
